\section{The limiting-threshold problem}\label{sec:intro}

As independent clauses are added to a random Boolean formula,
satisfiability changes from likely to unlikely. A narrow transition
window answers only part of the threshold question: its location might
still drift as the number of variables grows. We prove that the location
converges for every fixed clause size at least three.

A Boolean assignment gives each variable \(x_i\) a value in \(\{0,1\}\).
A literal is a variable or its negation; a clause is a disjunction of
literals, and a conjunctive normal form formula is a conjunction of
clauses. A formula is satisfiable if some assignment makes every clause
true. Random \(k\)-SAT asks how this event changes as independent clauses
are added to a formula on \(n\) variables.

Fix an integer \(k\ge3\). For \(n\ge k\), a \emph{proper \(k\)-clause}
uses \(k\) distinct variables chosen uniformly from \(x_1,\ldots,x_n\),
with independent fair signs. Thus each of the \(2^k\binom nk\) possible
clauses has the same probability. Let \(C_1,C_2,\ldots\) be independent
clauses with this law, and define
\[
 F_{n,m}=\bigwedge_{i=1}^{m}C_i,
 \qquad P_n(m)=\PP(F_{n,m}\in\sat),\qquad m=0,1,2,\ldots.
\]
Here \(\sat\) denotes satisfiability, and the empty conjunction is true.
Whole clauses are sampled with replacement: two positions may contain
the same clause, although a variable cannot occur twice within one
proper clause. The prefix coupling makes \(P_n\) nonincreasing and gives
\(P_n(0)=1\). The ratio \(c=m/n\) is the clause density.

\begin{theorem}\label{thm:main}
For every fixed integer \(k\ge3\), there is a constant
\(\alpha_k\in(0,\infty)\) such that, for independent proper uniformly
signed \(k\)-clauses sampled with replacement,
\[
 \lim_{n\to\infty}P_n(\lfloor cn\rfloor)
 =
 \begin{cases}
  1,&0\le c<\alpha_k,\\
  0,&c>\alpha_k.
 \end{cases}
\]
\end{theorem}

The conclusion is a strict-side threshold statement; it makes no
assertion at \(c=\alpha_k\). All estimates below may depend on the fixed
integer \(k\). The proof does not require the centers of the transition
windows to be known in advance. It constructs a finite-size center,
controls the fluctuations about it, and then proves that these centers
converge.

\subsection{Prior work and the role of the present argument}

Carenini~\cite[Theorem~1.1 and Corollary~1.2]{Carenini2026Polynomial}
proves an \(O_{k,\eta}(n^{1/2+1/k})\) clause window between probability
levels \(\eta\) and \(1-\eta\), for each fixed \(0<\eta<1/2\), and, using
the Abbe--Montanari criterion, establishes the satisfiability conjecture for
every fixed \(k\ge3\). We credit Carenini with priority for the resolution
of the satisfiability conjecture. The present paper gives an alternative
proof of the limiting-threshold conclusion.

Chv\'atal and Reed~\cite{CR} explicitly formulated the conjecture that
each fixed clause size greater than one has a limiting satisfiability
density. Friedgut's sharp-threshold theorem, with Bourgain's appendix
\cite[Theorem~1.3]{Friedgut}, established that random \(k\)-SAT has a
sharp transition around a sequence of densities. It left open whether
that sequence converges for each fixed \(k\). The classical two-clause
problem had already been settled by Chv\'atal and Reed
\cite[Theorems~3--4]{CR}, who also credit independent proofs by Goerdt
and Fernandez de la Vega. Goerdt's journal treatment is~\cite{Goerdt}.
The implication-graph mechanism is separate from the proof given here.

For larger arities, Achlioptas and Peres~\cite{AP} obtained leading-order
bounds \(2^k\log2-O(k)\), using weighted second moments. The
statistical-physics work of M\'ezard, Parisi, and Zecchina~\cite{MPZ}
and of Mertens, M\'ezard, and Zecchina~\cite{MMZ} developed threshold
predictions through the cavity method and one-step replica symmetry
breaking. In particular, the latter work predicted the large-\(k\)
expansion
\[
 2^k\log2-\frac{1+\log2}{2}+o_k(1),
 \qquad o_k(1)\longrightarrow0\quad(k\longrightarrow\infty).
\]
Coja-Oghlan and Panagiotou~\cite[Theorem~1.1]{COP} proved that the
liminf and limsup of the sharp-threshold sequence both lie within an
additive error tending to zero as \(k\to\infty\) of
\(2^k\log2-(1+\log2)/2\). These asymptotic-in-arity estimates are
distinct from convergence in \(n\) at a fixed arity. Ding, Sly, and
Sun~\cite{DSS} proved existence and identified the threshold with the
one-step replica symmetry breaking prediction for all sufficiently
large fixed \(k\). Theorem~\ref{thm:main} establishes existence for every
fixed \(k\ge3\), without identifying the value of \(\alpha_k\).

The first ingredient is a comparison between shorter clauses and
blocks of ordinary clauses. Friedgut's half-cube replacement lemma
\cite[Lemma~5.7]{Friedgut} gives a qualitative antecedent based on frozen
coordinates and sequential replacement. Carenini
\cite[Lemma~6.1 and Corollaries~6.2, 7.10--7.11]{Carenini} developed an
averaged clause-replacement method. The cited version proves an
\(O_k(n/\log n)\) window between any two fixed central probability
levels, with constants also depending on those levels, and distinguishes
window width from convergence of the location in Section~10.1.
Here the replacement is applied after erasing one variable. Summing the
resulting changes in satisfiability probabilities bounds the forced
fraction along the clause process. A deletion argument then converts
this integrated bound into polynomial concentration. We give the
replacement, forcing, and deletion proofs in full.

The second ingredient compares a whole system with independent
subsystems. Guerra and Toninelli~\cite{GT} used interpolation to
establish thermodynamic limits for mean-field spin glasses.
Franz and Leone~\cite{FL}, and Franz, Leone, and Toninelli~\cite{FLT},
developed interpolation for diluted systems. Their pressure and
free-energy arguments, with the stated even-arity restrictions, do not
supply the hard satisfiability-threshold conclusion used here.
Bayati, Gamarnik, and Tetali~\cite[Section~4, Proposition~2]{BGT}
established a frozen-variable satisfiability interpolation between
clauses on a whole variable set and clauses confined to its parts.
Our Poisson comparison is a continuous-parameter form of that method.
Its sign follows from convexity of \(x\mapsto x^k\), applied to the
proportions of globally forced variables in the parts. The comparison
is exact for an auxiliary clause model that permits repeated variables;
we prove the transfer back to proper clauses, including the constant
loss caused by repetitions.

Finally, errors that are summable over geometrically increasing sizes
cannot produce macroscopic drift. This principle parallels the
approximate-superadditivity analysis of Abbe and Montanari
\cite[arXiv version, Section~4.3, Lemma~8 and Remark~3]{AM}, developed
in their study of solution counts. Our comparison uses the minimum of
two normalized centers. We prove the corresponding deterministic
convergence lemma, in a form that requires the comparison only when
the two sizes are comparable.

\subsection{How the proof reaches the threshold}

We cap the last satisfiable index at a fixed multiple of \(n\), and
take its expectation divided by \(n\) as the finite-size center. The
first task is to show that this capped index has sublinear fluctuations.
In a satisfiable formula, a variable is \emph{forced} if every
satisfying assignment gives it the same value. The clauses that avoid
the variable form the untouched background. If the variable is forced,
deleting its literal from each incident clause produces shorter clauses
that eliminate every solution of this background. Replacing each
shortened clause by a block of ordinary clauses turns this event into
a decrease of the satisfiability probability. Summing over prefix
lengths telescopes the decreases. Section~\ref{sec:forcing} thereby
controls the total expected forced fraction along the process.

Section~\ref{sec:concentration} omits a single clause while keeping the
other time indices. The delay in losing satisfiability counts the
prefixes where that clause is pivotal. Its square counts pairs of such
prefixes. If the punctured formula survives from the earlier prefix to
the later one, every intervening clause must avoid killing the earlier
solution set. This bounds the pair probability in terms of the variables
already forced at the earlier time. The integrated forcing estimate
then yields the required variance bound. Section~\ref{sec:anchors}
places the normalized means in a fixed positive bounded interval, using
a clause-to-variable matching at low density and an assignment first
moment at high density.

Concentration alone still permits the centers to drift. In the auxiliary
Poisson model, interpolation says that satisfiability at a combined size
is at least the product of the probabilities in two independent parts.
Below either finite-size center, transfer from the proper model gives a
fixed positive probability; above the combined center, it gives a
probability tending to zero. These facts prevent the combined center
from lying much below both smaller centers. A comparison with one extra
variable handles adjacent sizes. Section~\ref{sec:convergence} makes
these steps precise and sums the errors along a balanced tree of nearly
equal blocks. Its concentration-to-convergence theorem is then applied
to the local variance and center bounds.

Section~\ref{sec:three-clause} uses a first moment at density six to
bound the means of the three-clause indices, and gives exact relations
between their cap conventions. It then specializes the forcing estimate and deduces
concentration at the smaller cap. The improved mean bound permits a
smaller transfer interval for either cap, sharpening the numerical
constants while keeping the roles of the cap and the interval distinct.
Section~\ref{sec:boundaries} explains the lower-arity cases and the
finite-size information supplied by the proof.
